\documentclass[11pt, a4paper]{article}

%% ─── Codificação e idioma ────────────────────────────────────────────────────
\usepackage[utf8]{inputenc}
\usepackage[T1]{fontenc}
\usepackage[portuguese]{babel}

%% ─── Tipografia ──────────────────────────────────────────────────────────────
\usepackage{lmodern}
\usepackage{microtype}

%% ─── Geometria de página ─────────────────────────────────────────────────────
\usepackage[top=2.5cm, bottom=2.5cm, left=3cm, right=3cm]{geometry}

%% ─── Matemática ──────────────────────────────────────────────────────────────
\usepackage{amsmath}
\usepackage{amssymb}

%% ─── Figuras e tabelas ───────────────────────────────────────────────────────
\usepackage{graphicx}
\usepackage{booktabs}
\usepackage{array}
\usepackage{caption}
\usepackage{xcolor}
\usepackage{float}

%% ─── Listas ──────────────────────────────────────────────────────────────────
\usepackage{enumitem}
\setlist[itemize]{noitemsep, topsep=4pt}
\setlist[enumerate]{noitemsep, topsep=4pt}

%% ─── Hipervínculos ───────────────────────────────────────────────────────────
\usepackage[hidelinks, colorlinks=true, linkcolor=blue!70!black,
            citecolor=blue!70!black, urlcolor=blue!70!black]{hyperref}

%% ─── Captions ────────────────────────────────────────────────────────────────
\captionsetup{font=small, labelfont=bf, labelsep=period}

%% ─── Comandos de conveniência ────────────────────────────────────────────────
\newcommand{\Rx}{\mathrm{Rx}}
\newcommand{\Tx}{\mathrm{Tx}}
\newcommand{\dxy}{\delta_{xy}}
\newcommand{\dz}{\delta_z}
\newcommand{\alphares}{\alpha_{\mathrm{res}}}

%% ─── Caminho comum das figuras geradas ───────────────────────────────────────
\newcommand{\figdir}{../simulations/io/plan_vuelo_multislab/analisis_comparativo}

%% ─── Metadados ───────────────────────────────────────────────────────────────
\title{\textbf{Sensibilidade do Otimizador de Trajetória do Receptor\\
       ao Peso Potência--Resolução e ao Tamanho da Grade de Alvos\\
       \large{Pipeline biestático multicamadas (multislab)}}}
\author{Nota técnica --- doutorado SAR}
\date{\today}

% =============================================================================
\begin{document}
% =============================================================================

\maketitle
\thispagestyle{empty}

\noindent\textbf{Contexto:} otimizador de posição do receptor (Rx) para o SAR
biestático helicoidal com propagação em solo multicamadas. O custo escalarizado
combina potência recebida e resolução volumétrica através do peso
$\alphares \in [0,1]$ (0 = só potência, 1 = só resolução). Este relatório resume
um estudo de sensibilidade desse peso e do tamanho da grade de alvos utilizada
para avaliar a qualidade da solução.
\vspace{6pt}

\noindent\textbf{Nota de revisão:} uma versão anterior deste relatório continha
um erro na rotina de registro de potência do otimizador
(\texttt{run\_plano\_de\_voo\_multislab.m}): a chamada que deveria medir a
potência física no ponto ótimo reutilizava por engano o mesmo $\alphares$ da
otimização, em vez de $\alphares=0$. Como consequência, a coluna de "potência"
reportada para $\alphares \neq 0$ não era potência, e sim uma reformulação de
$10^{-J}$ — o que produzia uma correlação $\alphares$--potência positiva,
fisicamente incoerente. O erro foi corrigido, os 12 cenários foram
reprocessados do zero e todos os números, figuras e conclusões abaixo refletem
os dados corrigidos.
\vspace{6pt}

\noindent\rule{\linewidth}{0.4pt}

% =============================================================================
\section{Objetivo do estudo}
% =============================================================================

O otimizador escolhe, para cada posição do transmissor (Tx) ao longo da órbita
helicoidal, a posição do receptor que minimiza o custo
\[
J(\Rx;\alphares) = -(1-\alphares)\,\log_{10} P_r(\Rx) + \alphares\,\log_{10}
\bigl[\dxy(\Rx)\cdot\dz(\Rx)\bigr],
\]
em que $P_r$ é a potência recebida e $\dxy$, $\dz$ são as resoluções lateral e
axial esperadas no centro da grade de alvos. Dois parâmetros de projeto afetam
diretamente essa otimização e não haviam sido caracterizados sistematicamente:

\begin{itemize}
  \item $\alphares$ --- o peso de compromisso entre potência e resolução;
  \item o tamanho da grade de alvos (extensão em $x$/$y$ da região de interesse
        no subsolo), que define sobre quantos pontos a resolução é avaliada e
        quão descentralizados do eixo da órbita eles podem estar.
\end{itemize}

O estudo varreu $\alphares \in \{0,\ 0.3,\ 0.6,\ 1\}$ e três tamanhos de grade
($30$, $60$ e $100$~m de lado), com $5\times5\times5 = 125$ alvos por cenário,
totalizando 12 execuções completas da otimização ao longo de toda a trajetória
do Tx.

% =============================================================================
\section{Resultados}
% =============================================================================

A Tabela~\ref{tab:resumo} apresenta, para cada combinação de grade e
$\alphares$, os valores médios ao longo da trajetória de: potência recebida,
resoluções lateral/axial e o produto $\dxy\cdot\dz$ (proxy do volume da célula
de resolução no plano horizontal).

\begin{table}[H]
  \centering
  \caption{Métricas médias por cenário (125 alvos, médias ao longo da trajetória do Tx).}
  \label{tab:resumo}
  \small
  \begin{tabular}{@{}ccrrrr@{}}
    \toprule
    \textbf{Grade (m)} & $\boldsymbol{\alphares}$ &
    \textbf{Potência (dBm)} & $\boldsymbol{\dxy}$\textbf{ (cm)} &
    $\boldsymbol{\dz}$\textbf{ (cm)} & $\boldsymbol{\dxy\!\cdot\!\dz}$\textbf{ (cm\textsuperscript{2})} \\
    \midrule
    30  & 0.0 & $-24.6$ & 1748.7 & 89.8 & $1.58\times10^{5}$ \\
    30  & 0.3 & $-26.2$ & 24.4   & 88.8 & 2169 \\
    30  & 0.6 & $-28.2$ & 17.5   & 93.2 & 1628 \\
    30  & 1.0 & $-31.0$ & 15.6   & 98.5 & 1540 \\
    \midrule
    60  & 0.0 & $-27.9$ & 1154.1 & 89.8 & $1.04\times10^{5}$ \\
    60  & 0.3 & $-29.2$ & 21.8   & 89.3 & 1943 \\
    60  & 0.6 & $-31.3$ & 17.2   & 93.9 & 1611 \\
    60  & 1.0 & $-33.8$ & 15.6   & 98.5 & 1540 \\
    \midrule
    100 & 0.0 & $-31.0$ & 2196.8 & 89.7 & $1.98\times10^{5}$ \\
    100 & 0.3 & $-32.4$ & 25.5   & 89.3 & 2275 \\
    100 & 0.6 & $-34.5$ & 17.1   & 94.3 & 1608 \\
    100 & 1.0 & $-36.8$ & 15.6   & 98.5 & 1540 \\
    \bottomrule
  \end{tabular}
\end{table}

\begin{figure}[H]
  \centering
  \includegraphics[width=0.85\linewidth]{\figdir/tradeoff_potencia_resolucion.png}
  \caption{Curva de compromisso potência--resolução. Cada ponto de uma mesma
  série (tamanho de grade) corresponde a um valor de $\alphares$; o eixo
  horizontal está em escala logarítmica.}
  \label{fig:tradeoff}
\end{figure}

\begin{figure}[H]
  \centering
  \includegraphics[width=0.95\linewidth]{\figdir/evolucion_grid_30m.png}
  \caption{Evolução ao longo da trajetória do Tx para a grade de 30~m: potência
  recebida (topo), produto de resolução $\dxy\cdot\dz$ em escala logarítmica
  (meio) e custo $J$ (base), comparando os quatro valores de $\alphares$.}
  \label{fig:evol30}
\end{figure}

\begin{figure}[H]
  \centering
  \includegraphics[width=0.95\linewidth]{\figdir/evolucion_grid_100m.png}
  \caption{Mesma análise da Fig.~\ref{fig:evol30}, para a grade de 100~m. O
  padrão qualitativo em função de $\alphares$ é o mesmo; a diferença está no
  nível absoluto de potência (ver Seção~\ref{sec:interpretacao}).}
  \label{fig:evol100}
\end{figure}

% =============================================================================
\section{Interpretação}
\label{sec:interpretacao}
% =============================================================================

\subsection{Efeito do peso $\alphares$}

O peso $\alphares$ controla um compromisso monotônico e fisicamente coerente:
a correlação entre $\alphares$ e a potência média é $-1.00$ nos três tamanhos
de grade (aumentar $\alphares$ \emph{sempre} custa potência, como esperado ao
despriorizá-la no custo), e a correlação entre $\alphares$ e a resolução
(produto $\dxy\cdot\dz$) é $-0.74$ (aumentar $\alphares$ tende a melhorar a
resolução). O ponto central da análise é o quão favorável é essa troca ao
longo do intervalo $[0,1]$ — e ela está longe de ser linear:

\begin{itemize}
  \item \textbf{De $\alphares=0$ para $\alphares=0.3$}: o custo em potência é
        pequeno ($1.4$--$1.6$~dB, consistente nos três tamanhos de grade),
        enquanto o produto $\dxy\dz$ cai entre $53\times$ e $87\times$
        (de $\sim\!10^5$--$2\times10^5$ para $\sim\!2\times10^3$
        cm\textsuperscript{2}, i.e. mais de $98\%$ de redução). Esta é, de
        longe, a transição de maior retorno: sair de $\alphares=0$ (que ignora
        completamente a resolução, deixando o Rx numa geometria com dezenas de
        metros de erro lateral) custa menos de $2$~dB de potência.
  \item \textbf{De $\alphares=0.3$ para $\alphares=0.6$}: o custo sobe para
        $\sim\!2$~dB, e o ganho de resolução cai para $17$--$29\%$ de redução
        adicional em $\dxy\dz$ — ainda vantajoso, mas com retorno claramente
        decrescente frente ao passo anterior.
  \item \textbf{De $\alphares=0.6$ para $\alphares=1.0$}: o custo em potência
        é o maior de todos ($2.3$--$2.9$~dB), para um ganho de resolução
        marginal de apenas $4$--$5\%$. Praticamente todo o benefício de
        resolução alcançável já foi obtido em $\alphares=0.6$.
  \item $\alphares=1$ converge para a mesma geometria relativa de resolução
        ($\dxy\approx15.6$~cm, $\dz\approx98.5$~cm) independentemente do
        tamanho da grade — esperado, pois quando o custo ignora totalmente a
        potência, a posição ótima do Rx é ditada apenas pelo ângulo biestático
        em relação ao centro da grade, não pela extensão da grade em si.
        Ainda assim, a \emph{potência física} nessa posição varia com o
        tamanho da grade ($-31.0 \to -33.8 \to -36.8$~dBm), porque os alvos
        mais afastados do centro continuam contribuindo menos à potência
        total mesmo quando a posição do Rx não foi escolhida em função deles.
\end{itemize}

\textbf{Recomendação prática:} a faixa $\alphares \in [0.3,\,0.6]$ é o ponto de
operação recomendado: captura mais de $98\%$ do ganho de resolução disponível
a um custo de potência de apenas $1.4$--$3.6$~dB frente a $\alphares=0$.
Empurrar $\alphares$ até $1.0$ custa outros $2$--$3$~dB de potência por um
ganho de resolução essencialmente nulo, e por isso não é recomendado como
padrão, exceto quando a resolução lateral for a única prioridade e houver
margem de SNR de sobra.

\subsection{Efeito do tamanho da grade de alvos}

Para $\alphares$ fixo, aumentar o tamanho da grade de alvos (30~$\to$~60~$\to$~100~m)
degrada monotonicamente a potência média recebida, com uma penalidade
notavelmente uniforme de $\sim\!3$~dB por salto de escala de grade,
\emph{independente do valor de $\alphares$ usado}:

\begin{itemize}
  \item $\alphares=0$: $-24.6 \to -27.9 \to -31.0$~dBm ($-3.3$, $-3.2$~dB).
  \item $\alphares=0.3$: $-26.2 \to -29.2 \to -32.4$~dBm ($-3.1$, $-3.1$~dB).
  \item $\alphares=0.6$: $-28.2 \to -31.3 \to -34.5$~dBm ($-3.1$, $-3.2$~dB).
  \item $\alphares=1.0$: $-31.0 \to -33.8 \to -36.8$~dBm ($-2.8$, $-3.0$~dB).
\end{itemize}

Notavelmente, a penalidade persiste mesmo em $\alphares=1$, quando a posição do
Rx é escolhida ignorando completamente a potência: mesmo que a geometria ótima
de resolução não dependa da grade, alvos mais espalhados continuam, em média,
mais distantes e pior iluminados por essa geometria fixa, reduzindo a potência
física recebida. Isso confirma que o efeito do tamanho de grade sobre a
potência é uma propriedade da própria distribuição espacial dos alvos, não um
artefato do peso de otimização.

Em compensação, o efeito da grade sobre a resolução é pequeno (comparar
colunas $\dxy\dz$ entre tamanhos de grade, para $\alphares$ fixo): a resolução
é dominada pela geometria biestática global (posição relativa Tx--Rx--centro
da grade) e não pela dispersão espacial dos alvos individuais.

\textbf{Implicação prática:} grades de alvos maiores (regiões de interesse mais
extensas) custam SNR de forma consistente, mas não custam resolução. Se a área
de interesse real for grande, convém compensar com potência de transmissão
adicional ou reprocessar por sub-regiões menores em vez de otimizar uma única
geometria de Rx para toda a área.

\subsection{Sobre o uso do custo $J$ como métrica de comparação}

O custo $J$ \emph{não} deve ser comparado diretamente entre corridas com
$\alphares$ distinto, pois sua própria definição muda com o peso (os valores de
$J$ na Tabela~\ref{tab:resumo}, se olhados fora de contexto, sugeririam
piora monotônica de $\alphares=0$ para $\alphares=1$, o que é apenas um
artefato da escala, não uma piora real de desempenho). As comparações de
desempenho neste relatório usam sempre as grandezas físicas diretamente
(potência em dBm, $\dxy$, $\dz$), como na Fig.~\ref{fig:tradeoff}.

% =============================================================================
\section{Conclusões}
% =============================================================================

\begin{enumerate}
  \item O peso $\alphares$ tem efeito monotônico e fisicamente coerente:
        aumentá-lo sempre custa potência (correlação $-1.00$) e sempre tende a
        melhorar a resolução (correlação $-0.74$), confirmando que o custo
        escalarizado está corretamente formulado. A troca é fortemente
        não linear: mais de $98\%$ do ganho de resolução disponível é obtido
        já em $\alphares=0.3$, a um custo de apenas $1.4$--$1.6$~dB de
        potência; valores acima de $0.6$ trocam potência adicional por uma
        melhora de resolução praticamente nula.
  \item O tamanho da grade de alvos afeta a potência recebida média de forma
        consistente e uniforme ($\sim$3~dB de perda por salto de escala de
        grade), \emph{para qualquer valor de $\alphares$, inclusive
        $\alphares=1$} — ao contrário do que uma versão anterior deste
        relatório (com o bug de potência ainda presente) sugeria. O impacto
        sobre a resolução esperada, em contraste, é pequeno.
  \item A superfície de compromisso potência--resolução (Fig.~\ref{fig:tradeoff})
        é a ferramenta correta para escolher $\alphares$ em função do requisito
        de projeto (SNR mínimo vs. resolução mínima), e não os valores brutos
        de $J$, que não são comparáveis entre diferentes $\alphares$.
  \item Para o cenário multicamadas avaliado, $\alphares \in [0.3,\,0.6]$ é
        recomendado como faixa de operação padrão — $0.3$ quando a potência
        recebida for a restrição mais crítica, $0.6$ quando a resolução
        lateral for prioritária e houver margem de SNR. $\alphares=1$ não é
        recomendado como padrão: seu custo adicional de potência
        ($2$--$3$~dB frente a $\alphares=0.6$) não se traduz em ganho
        perceptível de resolução.
\end{enumerate}

\end{document}
